\documentclass[12pt,a4paper]{article}
\usepackage[utf8]{inputenc}
\usepackage[italian]{babel}
\usepackage{amsmath,amssymb,amsfonts}
\usepackage{geometry}
\geometry{top=2cm,bottom=2cm,left=2cm,right=2cm}
\pagestyle{empty}

\begin{document}

\begin{center}
    {\large\bf Scheda 1. Principio dei Cassetti} \\
    \vspace{2mm}
    {\small Club di Matematica}
\end{center}

\vspace{5mm}

\begin{enumerate}
    \item Una scatola contiene 10 calzini neri e 10 calzini bianchi. Qual è il numero minimo di calzini che devi estrarre alla cieca per garantire di avere almeno un paio di calzini dello stesso colore?
    
    \item In una classe ci sono 30 studenti. Dimostra che ci sono almeno due studenti i cui cognomi iniziano con la stessa lettera.
    
    \item Dimostra che tra $n+1$ numeri interi qualsiasi, se ne possono scegliere due la cui differenza è divisibile per $n$.
\end{enumerate}

\end{document}
